\documentclass[10pt,a4paper]{amsart}
\usepackage[margin=19mm]{geometry}
\usepackage{amsmath,amssymb,mathtools}
\usepackage[scr=rsfs,cal=euler]{mathalfa}
\usepackage{xfrac}
\usepackage[unicode,hidelinks,bookmarks=false]{hyperref}
\newcommand{\ie}{i.e.\ }
\newcommand{\cf}{cf.\ }
\newcommand{\resp}{respectively\ }
\newcommand{\Subsection}{\subsection}
\providecommand{\nobreakdash}{\nobreak\hbox{-}}
\setlength{\emergencystretch}{2em}
\multlinegap0pt
\usepackage{bm}
\usepackage{bbm}
\usepackage{dsfont}
\usepackage{xspace}
\usepackage{cancel}
\usepackage{mathtools}
\usepackage{amsthm}
\makeatletter
\@ifundefined{corollary}{\newtheorem{corollary}{Corollary}}{}
\@ifundefined{theorem}{\newtheorem{theorem}{Theorem}}{}
\makeatother
\newcommand{\BB}{\mathbf{B}} % B-basis
\newcommand{\Comp}{\operatorname{Comp}}
\newcommand{\NN}{\mathbb{N}} % nonnegative integers
\newcommand{\VVq}{\mathbf{V}^{(q)}} % Dynkin operator
\newcommand{\comp}{\operatorname{comp}}
\newcommand{\kk}{\mathbf{k}} % our base ring
\title[A permutation based approach to the $q$-deformation of the D]{A permutation based approach to the $q$-deformation of the Dynkin Operator}
\author{Darij Grinberg, Ekaterina A. Vassilieva}
\date{Source: arXiv:2504.09709v2 (CC BY 4.0)}
\begin{document}
\maketitle

\noindent\emph{Source and licence.} A permutation based approach to the $q$-deformation of the Dynkin Operator. Version arXiv:2504.09709v2 (CC BY 4.0), distributed under the Creative Commons Attribution 4.0 International licence, as recorded in the arXiv licence metadata for this exact version and shown on the abstract page https://arxiv.org/abs/2504.09709v2. Full rights notice: licenses/arxiv-2504.09709-CC-BY-4.0.txt.\par\smallskip

\noindent\textit{Editorial scope.} Counted unit: the Corollary whose source label is \texttt{corollary.rank} in the article (source0.tex lines 630--650, source label \texttt{corollary.rank}), delivered together with the article's own complete original proof of it (source0.tex lines 651--713, one \texttt{proof} environment of 63 lines). No mathematics is written, repaired or re-derived: statement and proof are the article's own text line for line. Required context retained uncounted: thm.vnq (lines 495--506); corollary.eigen (lines 614--621). Every definition, display and label of those lines is retained unchanged. Omitted from this package and not redistributed: 1--494, 507--613, 622--629, 714--762. Nothing inside a retained excerpt is omitted or altered: every retained body is delivered exactly as the article's lines are written, so no omission receipt of any kind accompanies this package. Citations: the retained excerpts carry no citation command at all, so no bibliography entry of the article is transcribed and no reference is redistributed. Reference adaptation: none was needed and none was made; every reference inside the delivered unit resolves inside it, so no unresolved reference or citation is produced. Printed slips kept literal: no printed word, symbol, hypothesis or number of the counted statement or of its proof was altered. Layout and class: the article's own class is \texttt{FPSAC2025} with packages including tikz-cd, bm, relsize, comment, biblatex; the wrapper is the lane's pinned \texttt{amsart} editorial wrapper at 10pt on A4, which restates the theorem-like environments the retained text needs and adds the article's own macro definitions that the retained text needs (\texttt{\textbackslash BB}, \texttt{\textbackslash Comp}, \texttt{\textbackslash NN}, \texttt{\textbackslash VVq}, \texttt{\textbackslash comp}, \texttt{\textbackslash kk}), with extra wrapper packages (bm, bbm, dsfont, xspace, cancel, mathtools). The delivered numbering is the unit's own. Assets: no retained span contains a figure environment, a graphics-inclusion command, a PDF-page inclusion, a table or a TikZ picture; no image file is copied into this package, the package declares no asset dependency and no image file is redistributed with it. Licence: version arXiv:2504.09709v2 of this article is distributed under the Creative Commons Attribution 4.0 International licence, as recorded in the arXiv licence metadata for this exact version and shown on the abstract page https://arxiv.org/abs/2504.09709v2. A full rights notice is delivered at licenses/arxiv-2504.09709-CC-BY-4.0.txt. Characters: every retained excerpt is pure ASCII, so the delivered engine, XeLaTeX with its default Latin Modern text font, reports no missing character.
\begin{theorem}
\label{thm.vnq}
Let $n$ be a positive integer and $q \in \mathbf{k}$ be an invertible parameter.
% This does not work for $n = 0$.
For $I, K \subseteq [n-1]$, the coefficient in $\mathbf{B}_K$ of $\mathbf{V}^{(q)}_n\mathbf{B}_I$ is $0$ if $I \nsubseteq K$ and otherwise given by
\begin{equation}
\label{equation.thm.vnq}
[\mathbf{B}_K]\mathbf{V}^{(q)}_n\mathbf{B}_I
= (-1)^{|K|} q^{n-1} (1-q^{-1})^{|I|} \prod_{v \in \overline{ K|I}}[v]_{q^{-1}}.
\end{equation}
Here, for any integer $v \in \NN$ and any $r \in \kk$, we define the $r$-integer $[v]_r$ to be $r^0 + r^1 + \cdots + r^{v-1}$. (Note that $[v]_{q^{-1}} = q^{-v+1} [v]_q$.)
\end{theorem}

\begin{corollary}
\label{corollary.eigen}
Fix $n>0$ and $q\neq 1$. The eigenvalues of $\mathbf{V}_n^{(q)}$ are the elements of the family $(e_I)_{I \subseteq [n-1]}$ defined by
\begin{equation}
 (1-q)e_I = \prod_{v \in \comp(I)} (1-q^v).
\end{equation}
Note that if $q$ is a $p$-th root of unity with $p>1$ and $\comp(I)$ contains a multiple of $p$, then $e_I = 0$. 
\end{corollary}

\begin{corollary}
\label{corollary.rank}
Let $n>0$. Let $\mathbf{k}$ be a field. The dimension of the subspace $\mathbf{V}_n^{(q)}\Sigma_n$ of the descent algebra $\Sigma_n$ is $2^{n-1}$ when $q$ is not a root of unity.
If $q$ is a primitive $p$-th root of unity with $p>1$,
\begin{equation}
\dim\left(\mathbf{V}_n^{(q)}\Sigma_n\right) = s_n^{(p)},
\end{equation}
where $s_n^{(p)}$ is the $n$-th Fibonacci number of order $p$ defined by $s_0^{(p)} = 0$ and $s_n^{(p)} = 2^{n-1}$ for all $1 \leq n < p$ and the recurrence relation
%
\begin{equation*}
s^{(p)}_n = 
%\begin{cases}
%0\; & \mbox{ if } \; n=0,\\
%2^{n-1}\; & \mbox{ if } \; 1\leq n < p,\\
s^{(p)}_{n-1} + s^{(p)}_{n-2} + \cdots + s^{(p)}_{n-p} %\; & \mbox{ if } \; 
\qquad \text{for all } n \geq p.
%\\
%\end{cases}
\end{equation*}
%
\end{corollary}

\begin{proof}
If $q$ is not a root of unity,
then all the eigenvalues $e_I$ of $\VVq_n$ are nonzero (see Corollary~\ref{corollary.eigen}),
and thus $\VVq_n$ is invertible, so that its image is the whole $\Sigma_n$. Next, assume that $q$ is a primitive $p$-th root of unity with $p>1$.
We let $\Comp_n^{(p)}$ denote the set of the compositions of $n$ that do not contain a multiple of $p$.
We have $|\Comp_n^{(p)}| = s_n^{(p)}$ easily\footnote{%
If $p >n$, then no composition of $n$ contains a multiple of $p$, and thus $|\Comp_n^{(p)}| = 2^{n-1}$.
When $n\geq p$, we split compositions in $\Comp_n^{(p)}$ according to their last entry $m$.
If $m <p$, we get a bijection with the compositions in $\Comp_{n-m}^{(p)}$ by removing the last entry.
If $m>p$, replace $m$ by $m-p$ to get a composition in $\Comp_{n-p}^{(p)}$.
This last construction is also bijective, and the desired recurrence relation follows.} and it remains to show that $\dim\left(\mathbf{V}_n^{(q)}\Sigma_n\right) = |\Comp_n^{(p)}|$.
To prove this, we denote
$
a_{IK} := [\mathbf{B}_K]\mathbf{V}^{(q)}_n\mathbf{B}_I  \text{ for all } I, K \subseteq [n-1] 
$
and argue in two steps.

%\begin{itemize}
%\item[(i)]
\textit{Proof of $\dim\left(\mathbf{V}_n^{(q)}\Sigma_n\right) \geq |\Comp_n^{(p)}|$:}
Recall that the operator $u \mapsto \mathbf{V}_n^{(q)} u$ on $\Sigma_n$ is triangular (the $a_{IK}$ are zero unless $I\subseteq K$). As a result, the dimension of $\mathbf{V}_n^{(q)}\Sigma_n$  is at least as large as the number of non-zero entries on the diagonal, i.e. the number of non-zero $e_I=a_{II}=[\mathbf{B}_I]\mathbf{V}^{(q)}_n\mathbf{B}_I$ when $I$ ranges over all subsets of $[n-1]$.
%According to Corollary \ref{corollary.eigen}, all these $e_I$ are non-zero if $q$ is not a root of unity.
Since $q$ is a primitive $p$-th root of unity, Corollary \ref{corollary.eigen} shows that we have $e_I =0$ iff $\comp(I)$ contains at least one multiple of $p$, that is, iff $\comp(I) \notin \Comp_n^{(p)}$.
As a result, $\dim\left(\mathbf{V}_n^{(q)}\Sigma_n\right) \geq |\Comp_n^{(p)}|$.

%\item[(ii)] 
\textit{Proof of $\dim\left(\mathbf{V}_n^{(q)}\Sigma_n\right) \leq |\Comp_n^{(p)}|$:}
The next step is to show that $\dim\left(\mathbf{V}_n^{(q)}\Sigma_n\right) \leq |\Comp_n^{(p)}|$.
By the rank-nullity theorem, it will suffice to find $|\Comp_n \setminus \Comp_n^{(p)}|$ many linearly independent linear relations between the coefficients of each element of $\mathbf{V}_n^{(q)}\Sigma_n$.
Let $K \subseteq [n-1]$ be such that $\comp(K) = (\alpha_1, \alpha_2, \dots, \alpha_{|K|+1}) \in \Comp_n \setminus \Comp_n^{(p)}$.
Thus, some $\alpha_i$ is a multiple of $p$.
Let $m$ be maximal such that $\alpha_m$ is a multiple of $p$. 
\begin{itemize}
\item[(a)] If $m=|K|+1$ then
\begin{align}
\label{eq.linrel0q.1}
\forall f \in \mathbf{V}_n^{(q)}\Sigma_n, \ \ [\BB_K]f = 0.
\end{align}
Indeed, by linearity, it suffices to show that $a_{IK} = 0$ for each $I \subseteq [n-1]$. Either $I \nsubseteq K$ and $a_{IK} = 0$, or $I \subseteq K$ and (by construction) $\overline{K|I}$ contains $\alpha_{|K|+1}$ and thus $a_{IK}$ contains the factor $[\alpha_{|K|+1}]_{q^{-1}} = 0$ (see Theorem \ref{thm.vnq}). 
\item[(b)] Next assume $m < |K|+1$.
%The next part $\alpha_{m+1}$ is not a multiple of $p$.
Then, define the set $K' \subseteq [n-1]$ such that 
\[
\comp(K') = (\alpha_1, \dots, \alpha_{m-1}, \alpha_m + \alpha_{m+1}, \alpha_{m+2},\dots, \alpha_{|K|+1}).
\]
As a result $K' \subset K$. The composition $\comp(K')$ has exactly one part less than $\comp(K)$.
We fix any $I \subseteq [n-1]$, and aim to show that $a_{IK'} = - a_{IK}$.
Indeed, if $I \nsubseteq K$ and $I \nsubseteq K'$, then $a_{IK}=a_{IK'}=0$. If $I \subseteq K$ but $I \nsubseteq K'$, then $\alpha_m \in \overline{K|I}$ and $a_{IK} = 0$ (having a factor of $[\alpha_m]_{q^{-1}}$). But $a_{IK'} =0$ as well, since $I \nsubseteq K'$.
Finally, if $I \subseteq K' \subseteq K$, then either $\alpha_{m+1} \notin \overline{K|I}$ and the only difference between $a_{IK}$ and $a_{IK'}$ is the coefficient $(-1)^{|K|} = - (-1)^{|K'|}$; or $\alpha_{m+1} \in \overline{K|I}$ and an additional difference between coefficients $a_{IK}$ and $a_{IK'}$ is that the factor $[\alpha_{m+1}]_{q^{-1}}$ in $a_{IK}$ is replaced by $[\alpha_{m}+\alpha_{m+1}]_{q^{-1}}$ in $a_{IK'}$. But as $q$ is a $p$-th root of unity and $\alpha_{m}$ is a multiple of $p$,
we have $[\alpha_{m}+\alpha_{m+1}]_{q^{-1}} = [\alpha_{m+1}]_{q^{-1}}$
and therefore $a_{IK'} = - a_{IK}$ for all $I$. By linearity, we conclude that
\begin{align}
\label{eq.linrel0q.2}
\forall f \in \mathbf{V}_n^{(q)}\Sigma_n, \ [\BB_{K'}]f = -[\BB_K]f.
\end{align}
\end{itemize}
Thus, for each $K \subseteq [n-1]$ with $\comp(K) \in \Comp_n \setminus \Comp_n^{(p)}$, we have found a linear relation -- either \eqref{eq.linrel0q.1} or \eqref{eq.linrel0q.2} -- that holds for the coefficients of all $f \in \mathbf{V}_n^{(q)}\Sigma_n$.
These relations are linearly independent due to distinct ``leading terms''.
Hence, $\dim\left(\mathbf{V}_n^{(q)}\Sigma_n\right) \leq |\Comp_n^{(p)}|$ is proved,
so that $\dim\left(\mathbf{V}_n^{(q)}\Sigma_n\right) = |\Comp_n^{(p)}|$ follows.
%\qedhere
%\end{itemize}
\end{proof}

\end{document}
